% Part: first-order-logic
% Chapter: syntax-and-semantics
% Section: terms-formulas

\documentclass[../../../include/open-logic-section]{subfiles}

\begin{document}

\olfileid{fol}{syn}{frm}

\olsection{पदे आणि \printtoken{P}{formula}}

प्रथम-क्रम भाषा~$\Lang L$ दिली की, $\Lang L$ च्या मूलभूत
शब्दसंग्रहापासून तयार होणारी व्यंजके आपण परिभाषित करू शकतो. त्यांत
विशेषतः \emph{पदे} आणि \emph{!!{formula}s} यांचा समावेश होतो.

\begin{defn}[पदे]
\ollabel{defn:terms}
$\Lang L$ मधील \emph{पदांचा} संच~$\Trm[L]$ पुढीलप्रमाणे विगमनाने
परिभाषित केला जातो:
\begin{enumerate}
\item प्रत्येक !!{variable} हे पद आहे.
\item $\Lang L$ मधील प्रत्येक !!{constant} हे पद आहे.
\item $f$ हे $n$-स्थानी !!{function} आणि $t_1$, \dots, $t_n$ ही
  पदे असतील, तर $\Atom{f}{t_1, \ldots, t_n}$ हे पद आहे.
\tagitem{limitClause}{
  इतर काहीही पद नाही.}{}
\end{enumerate}
ज्या पदात एकही !!{variable}s नाही ते \emph{बंद पद} होय.
\end{defn}

\begin{explain}
भाषेचे विनिर्देशन करताना !!{constant}s ही पदांपासून वेगळी चिन्हांची
श्रेणी म्हणून दिसतात; पण त्याऐवजी ती शून्यस्थानी !!{function}s मध्ये
समाविष्ट करता आली असती. मग पदांच्या व्याख्येतील दुसरी बाब आवश्यक
राहिली नसती. $n = 0$ असेल तेव्हा $\Atom{f}{t_1, \ldots, t_n}$ याचा
अर्थ केवळ स्वतंत्र $f$ असा घ्यावा लागेल.
\end{explain}

\begin{defn}[सूत्रे]
\ollabel{defn:formulas}
$\Lang L$ या भाषेतील \emph{!!{formula}s} चा संच~$\Frm[L]$
पुढीलप्रमाणे विगमनाने परिभाषित केला जातो:
\begin{enumerate}
\tagitem{prvFalse}{$\lfalse$ हे आण्विक !!{formula} आहे.}{}

\tagitem{prvTrue}{$\ltrue$ हे आण्विक !!{formula} आहे.}{}

\item $R$ हे $\Lang L$ मधील $n$-स्थानी !!{predicate} आणि $t_1$, \dots,
  $t_n$ ही $\Lang L$ मधील पदे असतील, तर $\Atom{R}{t_1,\ldots, t_n}$ हे
  आण्विक !!{formula} आहे.

\item $t_1$ आणि $t_2$ ही $\Lang L$ मधील पदे असतील, तर $\Atom{\eq}{t_1, t_2}$
  हे आण्विक !!{formula} आहे.

\tagitem{prvNot}{$!A$ हे !!a{formula} असेल, तर $\lnot !A$ हे
  !!a{formula} आहे.}{}

\tagitem{prvAnd}{$!A$ आणि $!B$ ही !!{formula}s असतील, तर $(!A \land
  !B)$ हे !!a{formula} आहे.}{}

\tagitem{prvOr}{$!A$ आणि $!B$ ही !!{formula}s असतील, तर $(!A \lor !B)$
  हे !!a{formula} आहे.}{}

\tagitem{prvIf}{$!A$ आणि $!B$ ही !!{formula}s असतील, तर $(!A \lif !B)$
  हे !!a{formula} आहे.}{}

\tagitem{prvIff}{$!A$ आणि $!B$ ही !!{formula}s असतील, तर $(!A \liff !B)$
  हे !!a{formula} आहे.}{}

\tagitem{prvAll}{$!A$ हे !!a{formula} आणि $x$ हे !!a{variable} असेल,
  तर $\lforall[x][!A]$ हे !!a{formula} आहे.}{}

\tagitem{prvEx}{$!A$ हे !!a{formula} आणि $x$ हे !!a{variable} असेल,
  तर $\lexists[x][!A]$ हे !!a{formula} आहे.}{}

\tagitem{limitClause}{इतर काहीही !!a{formula} नाही.}{}
\end{enumerate}
\end{defn}

\begin{explain}
पदांच्या संचाची आणि !!{formula}s च्या संचाची व्याख्या या
\emph{विगमनाधारित व्याख्या} आहेत. मूलतः, आपण !!{formula}s चा संच
अनंतपणे अनेक टप्प्यांत तयार करतो. आरंभीच्या टप्प्यात सर्व आण्विक सूत्रे
ही सूत्रे आहेत असे घोषित करतो; हे व्याख्येतील पहिल्या काही बाबींशी,
म्हणजे
\iftag{prvTrue}{$\ltrue$, }{}%
\iftag{prvFalse}{$\lfalse$, }{}%
$\Atom{R}{t_1,\dots,t_n}$ आणि $\Atom{\eq}{t_1,t_2}$ यांच्या बाबींशी जुळते.
म्हणून ``आण्विक !!{formula}'' म्हणजे या रूपांपैकी कोणतेही !!{formula}.

व्याख्येतील इतर बाबी आधीच तयार केलेल्या !!{formula}s पासून नवी
!!{formula}s तयार करण्याचे नियम देतात. दुसऱ्या टप्प्यात या नियमांनी
आण्विक !!{formula}s पासून !!{formula}s तयार करता येतात. तिसऱ्या
टप्प्यात आण्विक सूत्रे आणि दुसऱ्या टप्प्यात मिळालेली सूत्रे यांपासून
नवी सूत्रे तयार करतो, आणि ही प्रक्रिया पुढे चालू राहते. अशा एखाद्या
टप्प्यात शेवटी तयार होणारी प्रत्येक गोष्ट !!{formula} असते; इतर
काहीही सूत्र नसते.
\end{explain}

संकेताप्रमाणे, आपण $\eq$ हे चिन्ह त्याच्या युक्तिवादांच्या मध्ये लिहून
कंस वगळतो: $\eq[t_1][t_2]$ हा $\Atom{\eq}{t_1,t_2}$ चा संक्षेप आहे.
तसेच, $\lnot \Atom{\eq}{t_1,t_2}$ चा संक्षेप $\eq/[t_1][t_2]$ असा
लिहितो. द्विस्थानी संयोजक~$\ast$ वापरून $!B$, $!C$ पासून तयार केलेले
सूत्र $(!B \ast !C)$ लिहिताना आपण अनेकदा सर्वांत बाहेरची कंसांची
जोडी वगळून केवळ~$!B \ast !C$ असे लिहू.

\begin{intro}
काही तर्कशास्त्राच्या ग्रंथांत
\iftag{prvEx}{$\lexists[x][!A]$ }{}%
\iftag{notprvEx,notprvAll}{}{आणि }%
\iftag{prvAll}{$\lforall[x][!A]$ }{}%
यांची गणना
\iftag{notprvEx,notprvAll}{!!a{formula}}{!!{formula}s} म्हणून करण्यासाठी
!!{variable}~$x$ हे~$!A$ मध्ये आलेले असणे आवश्यक मानतात. ही अट न
घातल्याने कोणतीही अडचण येत नाही आणि काम अधिक सोपे होते.
\end{intro}

\begin{tagblock}{defNot,defOr,defAnd,defIf,defIff,defTrue,defFalse,defEx,defAll}
\begin{defn}
परिभाषित कारके वापरून तयार केलेली सूत्रे पुढीलप्रमाणे समजायची:

\begin{tagenumerate}{defTrue,defFalse,defNot,defOr,defAnd,defIf,defIff,defEx,defAll}

\tagitem{defTrue}{$\ltrue$ हा
  \iftag{prvFalse}{$\lnot\lfalse$}{$(!A \lor \lnot !A)$ (येथे~$!A$
    हे एखादे निश्चित आण्विक !!{formula} आहे)} चा संक्षेप आहे.}{}

\tagitem{defFalse}{$\lfalse$ हा
  \iftag{prvTrue}{$\lnot\ltrue$}{$(!A \land \lnot !A)$ (येथे~$!A$
    हे एखादे निश्चित आण्विक !!{formula} आहे)} चा संक्षेप आहे.}{}

\tagitem{defNot}{$\lnot !A$ हा $!A \lif \lfalse$ चा संक्षेप आहे.}{}

\tagitem{defOr}{$!A \lor !B$ हा
  \iftag{prvAnd}{$\lnot(\lnot !A \land \lnot !B)$}{$\lnot !A \lif
    !B$} चा संक्षेप आहे.}{}

\tagitem{defAnd}{$!A \land !B$ हा
  \iftag{prvOr}{$\lnot(\lnot !A \lor \lnot !B)$}{$\lnot (!A \lif
    \lnot !B)$} चा संक्षेप आहे.}{}

\tagitem{defIf}{$!A \lif !B$ हा
  \iftag{prvOr}{$\lnot !A \lor !B$}{$\lnot (!A \land \lnot !B)$} चा संक्षेप आहे.}{}
%READERNOTE{OLFOL-005: गोठवलेल्या इंग्रजीतील पहिल्या संक्षेपात जुळणारा आरंभीचा कंस नसतानाही शेवटचा कंस आहे. मराठीत जुळत नसलेला कंस वगळला असून इंग्रजी बाइट्स बदललेले नाहीत.}

\tagitem{defIff}{$!A \liff !B$ हा $(!A \lif !B) \land (!B
  \lif !A)$ चा संक्षेप आहे.}{}

\tagitem{defAll}{$\lforall[x][!A]$ हा $\lnot\lexists[x][\lnot !A]$ चा संक्षेप आहे.}{}

\tagitem{defEx}{$\lexists[x][!A]$ हा $\lnot\lforall[x][\lnot !A]$ चा संक्षेप आहे.}{}
\end{tagenumerate}
\end{defn}
\end{tagblock}

विशिष्ट उपयोगासाठीच्या भाषेत काम करताना आपण अनेकदा द्विस्थानी
!!{predicate}s आणि !!{function}s त्यांच्या संबंधित पदांच्या मध्ये
लिहितो; उदा., अंकगणिताच्या भाषेत $t_1 < t_2$ आणि $(t_1 + t_2)$, तर
संचसिद्धान्ताच्या भाषेत $t_1 \in t_2$. अंकगणिताच्या भाषेतील उत्तराधिकारी
फलन तर संकेताप्रमाणे त्याच्या युक्तिवादाच्या \emph{नंतर} लिहिले
जाते:~$t'$. मात्र, औपचारिकदृष्ट्या हे अनुक्रमे
$\Atom{\Obj A^2_0}{t_1, t_2}$, $\Obj f^2_0(t_1, t_2)$,
$\Atom{\Obj A^2_0}{t_1, t_2}$ आणि $f^1_0(t)$ यांचे केवळ रूढ संक्षेप आहेत.

\begin{defn}[विन्यासात्मक अनन्यता]
$\ident$ हे चिन्ह चिन्हमालांमधील विन्यासात्मक अनन्यता व्यक्त करते;
म्हणजे $!A \ident !B$ तेव्हा आणि केवळ तेव्हाच, जेव्हा $!A$ आणि $!B$
या समान लांबीच्या चिन्हमाला असतात आणि प्रत्येक स्थानी त्यांच्यात तेच
चिन्ह असते.
\end{defn}

$\ident$ या चिन्हाच्या दोन्ही बाजूंना जोडणीने मिळालेल्या चिन्हमाला
येऊ शकतात. उदाहरणार्थ, $!A \ident (!B \lor !C)$ याचा अर्थ असा:
आरंभीचा कंस, चिन्हमाला~$!B$, $\lor$ हे चिन्ह, चिन्हमाला~$!C$ आणि
शेवटचा कंस या क्रमाने जोडून मिळणारी चिन्हमाला आणि चिन्हमाला~$!A$
या एकच आहेत. असे असेल, तर $!A$ चे पहिले चिन्ह आरंभीचा कंस आहे;
$!A$ मध्ये दुसऱ्या चिन्हापासून सुरू होणारी $!B$ ही उपचिन्हमाला आहे;
तिच्यानंतर $\lor$ येते; इत्यादी बाबी आपल्याला माहीत होतात.

पदे आणि !!{formula}s ही मूलभूत घटकांपासून विगमनाधारित व्याख्यांनी तयार
होत असल्यामुळे, त्यांच्याविषयीच्या गोष्टी सिद्ध करण्यासाठी पुढील विगमन
तत्त्वे वापरता येतात.

\begin{lem}[\emph{पदांवरील विगमनाचे तत्त्व}]
    \ollabel{lem:trmind}
    $\Lang L$ ही प्रथम-क्रम भाषा असू द्या.
    एखादा गुणधर्म~$P$ असा असेल की
    %
    \begin{enumerate}
        \item तो प्रत्येक !!{variable}~$v$ ला लागू होतो,
        %
        \item तो $\Lang L$ मधील प्रत्येक !!{constant}~$a$ ला लागू होतो, आणि
        %
        \item तो $t_1$,~\dots, $t_n$ यांना लागू होत असेल आणि $f$ हे
        $\Lang L$ मधील $n$-स्थानी !!{function} असेल, तर तो
        $f(t_1,\dotsc,t_n)$ लाही लागू होतो
    \end{enumerate}
    (येथे $t_1$,~\dots, $t_n$ ही $\Lang{L}$ मधील पदे आहेत असे गृहीत धरले आहे),
    तर $P$ हा $\Trm[L]$ मधील प्रत्येक पदाला लागू होतो.
\end{lem}

\begin{prob}
    \olref[fol][syn][frm]{lem:trmind} सिद्ध करा.
\end{prob}

\begin{lem}[\emph{!!{formula}s वरील विगमनाचे तत्त्व}]
    \ollabel{thm:frmind}
    $\Lang L$ ही प्रथम-क्रम भाषा असू द्या.
    एखादा गुणधर्म~$P$ सर्व आण्विक !!{formula}s ला लागू होत असेल आणि तो असा असेल की
    %
    \begin{enumerate}
        \tagitem{prvNot}{तो $!A$ ला लागू होत असेल, तर तो $\lnot !A$ ला
            लागू होतो;}{}
        \tagitem{prvAnd}{तो $!A$ आणि~$!B$ ला लागू होत असेल, तर तो
            $(!A \land !B)$ ला लागू होतो;}{}
        \tagitem{prvOr}{तो $!A$ आणि~$!B$ ला लागू होत असेल, तर तो
            $(!A \lor !B)$ ला लागू होतो;}{}
        \tagitem{prvIf}{तो $!A$ आणि~$!B$ ला लागू होत असेल, तर तो
            $(!A \lif !B)$ ला लागू होतो;}{}
        \tagitem{prvIff}{तो $!A$ आणि~$!B$ ला लागू होत असेल, तर तो
            $(!A \liff !B)$ ला लागू होतो;}{}
        \tagitem{prvEx}{तो $!A$ ला लागू होत असेल, तर तो
            $\lexists[x][!A]$ ला लागू होतो;}{}
        \tagitem{prvAll}{तो $!A$ ला लागू होत असेल, तर तो
            $\lforall[x][!A]$ ला लागू होतो;}{}
  \end{enumerate}
  (येथे $!A$ आणि $!B$ ही $\Lang{L}$ मधील !!{formula}s आहेत असे गृहीत धरले आहे),
  तर $P$ हा $\Frm[L]$ मधील सर्व सूत्रांना लागू होतो.
\end{lem}

\end{document}
